\section{2018 Analysis & Differential Equations}

\subsection{Individual}

\begin{exercise}
Let $(M_n,\|\cdot\|)$ be the Banach space formed by $n\times n$ real matrices, equipped with Hilbert-Schmidt norm. Prove:

\begin{itemize}
\item[(1)] If $\gamma:\mathbb{R}\to M_n$ is a $C^1$-function with $\gamma(0)=\mathbb{I}$ (the identity matrix) and $\dot{\gamma}(0)=A\in M_n$, then for $\forall t\in\mathbb{R}$, the sequence $\{\gamma^n(t/n)\}_{n=1}^{\infty}$ converges to $\exp(tA)$. In particular, for $\forall A,B\in M_n$, $e^{t(A+B)} = \lim_{n\to\infty}(e^{\frac{t}{n}A}e^{\frac{t}{n}B})^n$.
\item[(2)] For $\forall A,B\in M_n$, if $e^{t(A+B)} = e^{tA}e^{tB}$, $\forall t\in\mathbb{R}$, then $[A,B]:=AB-BA=0$.
\end{itemize}
\end{exercise}

\begin{solution}
\textbf{Pedagogical Motivation:} This exercise explores the fundamental relationship between a Lie group (the group of invertible matrices) and its Lie algebra (the space of matrices with the commutator bracket). Part (1) is a generalized version of the \textit{Lie Product Formula} (or Trotter Product Formula), which shows how the exponential of a sum can be recovered from the exponentials of individual components, even when they don't commute. Part (2) demonstrates that the identity $e^{A+B} = e^A e^B$ holds globally for all $t$ if and only if the operators commute, which is a key property of the exponential map.

(1) Since $\gamma$ is a $C^1$-function, we can use the Taylor expansion of $\gamma$ at $t=0$:
\[ \gamma(s) = \gamma(0) + \dot{\gamma}(0)s + o(s) = \mathbb{I} + As + o(s) \text{ as } s \to 0. \]
For a fixed $t \in \mathbb{R}$, letting $s = t/n$, we have:
\[ \gamma(t/n) = \mathbb{I} + \frac{tA}{n} + o(1/n) \text{ as } n \to \infty. \]
We know that for any $X \in M_n$, $( \mathbb{I} + \frac{X}{n} + o(1/n) )^n \to \exp(X)$ as $n \to \infty$. This follows from the fact that:
\[ n \ln(\mathbb{I} + \frac{X}{n} + o(1/n)) = n (\frac{X}{n} + o(1/n)) = X + o(1) \to X. \]
Taking the exponential yields the result. Thus, $\gamma^n(t/n) \to \exp(tA)$.
For the particular case, let $\gamma(t) = e^{tA} e^{tB}$. Then $\gamma(0) = \mathbb{I} \cdot \mathbb{I} = \mathbb{I}$. The derivative is:
\[ \dot{\gamma}(t) = \frac{d}{dt}(e^{tA}) e^{tB} + e^{tA} \frac{d}{dt}(e^{tB}) = A e^{tA} e^{tB} + e^{tA} B e^{tB}. \]
Thus $\dot{\gamma}(0) = A + B$. Applying the first part, we get:
\[ \lim_{n\to\infty} (e^{\frac{t}{n}A} e^{\frac{t}{n}B})^n = \exp(t(A+B)). \]

(2) Suppose $e^{t(A+B)} = e^{tA} e^{tB}$ for all $t \in \mathbb{R}$. Differentiating both sides with respect to $t$:
\[ (A+B) e^{t(A+B)} = A e^{tA} e^{tB} + e^{tA} B e^{tB}. \]
Substituting $e^{t(A+B)} = e^{tA} e^{tB}$ into the left side:
\[ (A+B) e^{tA} e^{tB} = A e^{tA} e^{tB} + e^{tA} B e^{tB}. \]
Subtracting $A e^{tA} e^{tB}$ from both sides gives:
\[ B e^{tA} e^{tB} = e^{tA} B e^{tB}. \]
Multiplying by $(e^{tB})^{-1} = e^{-tB}$ on the right, we obtain:
\[ B e^{tA} = e^{tA} B \text{ for all } t \in \mathbb{R}. \]
Differentiating this equation with respect to $t$ at $t=0$:
\[ B (A e^{0A}) = (A e^{0A}) B \implies BA = AB. \]
Thus, $[A,B] = AB - BA = 0$.
\end{solution}

\begin{exercise}
Let a second-order ordinary differential equation be given:
\begin{equation} \label{ODE_star} \tag{$*$}
\frac{d^2 x}{dt^2} = f(t,x,\frac{d}{dt}x),
\end{equation}
where $f(t,x,y)$ is a $C^1$ function on $[0,1]\times\mathbb{R}^2$. Let $\varphi$, $t\in[0,1]$ be a solution of the second-order such that $\varphi(0)=a$, $\varphi(1)=b$ where $a,b$ are given real numbers. Suppose $\frac{\partial f}{\partial x}(t,x,y) > 0$ for $(t,x,y)\in[0,1]\times\mathbb{R}^2$. Prove: if $|\beta-b|$ is sufficiently small, $\beta\in\mathbb{R}$, then there exists a solution $\psi$ of \eqref{ODE_star} such that $\psi(0)=a$, $\psi(1)=\beta$.
\end{exercise}

\begin{solution}
\textbf{Pedagogical Motivation:} This problem is a classic application of the \textit{Shooting Method} in the theory of boundary value problems (BVPs). It reduces the question of finding a BVP solution to an Initial Value Problem (IVP) depending on a parameter (the initial slope). The core analytic tool is the \textit{Inverse Function Theorem}, and the condition $\partial f / \partial x > 0$ ensures that the solution at the endpoint depends monotonically on the initial slope, preventing "focal points" or "conjugate points" in the interval $[0,1]$.

Consider the initial value problem for the ODE \eqref{ODE_star}:
\[ x''(t) = f(t, x(t), x'(t)), \quad x(0) = a, \quad x'(0) = \alpha. \]
For each $\alpha \in \mathbb{R}$, let $x(t, \alpha)$ be the unique solution to this initial value problem. By the theorem of smoothness with respect to initial data, $x(t, \alpha)$ is $C^1$ in $\alpha$. We are given that for some $\alpha_0$, $\varphi(t) = x(t, \alpha_0)$ satisfies $x(0, \alpha_0) = a$ and $x(1, \alpha_0) = b$.
To show that there exists a solution $\psi(t) = x(t, \alpha)$ such that $x(1, \alpha) = \beta$ for $\beta$ near $b$, it suffices by the Inverse Function Theorem to show that $\frac{\partial x}{\partial \alpha}(1, \alpha_0) \neq 0$.
Let $z(t) = \frac{\partial x}{\partial \alpha}(t, \alpha_0)$. Differentiating the ODE and initial conditions with respect to $\alpha$ (the variational equation):
\[ z''(t) = \frac{\partial f}{\partial x}(t, \varphi, \varphi') z(t) + \frac{\partial f}{\partial y}(t, \varphi, \varphi') z'(t), \quad z(0) = 0, \quad z'(0) = 1. \]
Let $Q(t) = \frac{\partial f}{\partial x}(t, \varphi(t), \varphi'(t)) > 0$ and $P(t) = \frac{\partial f}{\partial y}(t, \varphi(t), \varphi'(t))$. The equation for $z$ is:
\[ z'' - P(t) z' - Q(t) z = 0, \quad z(0) = 0, \quad z'(0) = 1. \]
Suppose $z(1) = 0$. Since $z(0) = 0$ and $z(1) = 0$, and $z$ is not identically zero (as $z'(0)=1$), $z$ must have a positive maximum or a negative minimum in $(0, 1)$ by Rolle's Theorem or the Extreme Value Theorem.
If $z$ has a positive maximum at $t_0 \in (0, 1)$, then $z(t_0) > 0$, $z'(t_0) = 0$, and $z''(t_0) \leq 0$. However, from the ODE:
\[ z''(t_0) = P(t_0) z'(t_0) + Q(t_0) z(t_0) = Q(t_0) z(t_0) > 0, \]
which is a contradiction. Similarly, if $z$ has a negative minimum, $z(t_0) < 0$, $z'(t_0) = 0$, $z''(t_0) \geq 0$, but the ODE gives $z''(t_0) < 0$.
Thus, $z(t)$ cannot be zero on $(0, 1]$. In particular, $z(1) = \frac{\partial x}{\partial \alpha}(1, \alpha_0) \neq 0$.
By the Implicit Function Theorem, the map $\alpha \mapsto x(1, \alpha)$ is a local diffeomorphism near $\alpha_0$. Thus, for $\beta$ sufficiently close to $b$, there exists an initial slope $\alpha$ such that $x(1, \alpha) = \beta$.
\end{solution}

\begin{exercise}
A function $\varphi:\mathbb{R}^n\to\mathbb{C}$ is called positive definite if for all $k\in\mathbb{N}$, $y_j\in\mathbb{R}^n$, $c_j\in\mathbb{C}$, $j=1,\dots,k$, one has $\sum_{i,j=1}^k c_i\bar{c}_j\varphi(y_i-y_j) \geq 0$. If $\varphi$ is a measurable positive definite function on $\mathbb{R}^n$. Prove:

\begin{itemize}
\item[(1)] $\varphi(-y) = \overline{\varphi(y)}$ and $|\varphi(y)|\leq\varphi(0)$.
\item[(2)] For every Lebesgue integrable nonnegative function $f$ on $\mathbb{R}^n$, one has
\[
\int_{\mathbb{R}^n}\int_{\mathbb{R}^n}\varphi(x-y)f(x)f(y)\,dx\,dy \geq 0.
\]
\item[(3)] If the function $f$ is also even, then
\[
\int_{\mathbb{R}^n}\varphi(x)f*f(x)\,dx \geq 0,
\]
where $*$ stands for convolution.
\item[(4)] For all $\alpha>0$, we have
\[
\int_{\mathbb{R}^n}\varphi(x)\exp(-\alpha|x|^2)\,dx \geq 0.
\]
\end{itemize}
\end{exercise}

\begin{solution}
\textbf{Pedagogical Motivation:} Positive definite functions are central to harmonic analysis and probability theory. \textit{Bochner's Theorem} states that every continuous positive definite function is the Fourier transform of a finite positive measure. This exercise guides us through the elementary properties of these functions and their behavior under convolution and integration, which eventually leads to the intuition behind Bochner's theorem.

(1) For $k=1$, $y_1=0$, we have $|c_1|^2 \varphi(0) \geq 0$, so $\varphi(0) \geq 0$.
For $k=2$, let $y_1=y, y_2=0$. The matrix $A = \begin{pmatrix} \varphi(0) & \varphi(y) \\ \varphi(-y) & \varphi(0) \end{pmatrix}$ must be positive semi-definite. A positive semi-definite matrix must be Hermitian, so $\varphi(-y) = \overline{\varphi(y)}$. Also, the determinant must be non-negative: $\varphi(0)^2 - \varphi(y)\varphi(-y) \geq 0$, which means $\varphi(0)^2 - |\varphi(y)|^2 \geq 0$, hence $|\varphi(y)| \leq \varphi(0)$.

(2) The definition $\sum_{i,j} c_i \bar{c}_j \varphi(y_i-y_j) \geq 0$ implies that for any simple function $g(x) = \sum c_i \chi_{E_i}(x)$, the integral $\int \int \varphi(x-y) g(x) \overline{g(y)} dx dy \geq 0$. Since $f$ is integrable and non-negative, we can approximate $f$ by simple functions in $L^1$. Since $\varphi$ is bounded by $\varphi(0)$, the integral converges, and by the continuity of the integral with respect to $L^1$ convergence, $\int \int \varphi(x-y) f(x) f(y) dx dy \geq 0$ holds.

(3) If $f$ is even, then $f(x) = f(-x)$. The convolution is $(f*f)(x) = \int f(x-y)f(y) dy$.
By \textbf{Bochner's Theorem}, a measurable positive definite function is the Fourier transform of a positive finite measure $\mu$: $\varphi(x) = \int_{\mathbb{R}^n} e^{ix \cdot \xi} d\mu(\xi)$.
Then $\int \varphi(x) (f*f)(x) dx = \int (\int e^{ix \cdot \xi} d\mu(\xi)) (f*f)(x) dx = \int \widehat{f*f}(\xi) d\mu(\xi)$.
Since $\widehat{f*f}(\xi) = \hat{f}(\xi)^2$ and $f$ is real and even, $\hat{f}(\xi)$ is real and even. Thus $\hat{f}(\xi)^2 \geq 0$.
Integrating a non-negative function against a positive measure gives a non-negative result.

(4) Let $G(x) = \exp(-\alpha |x|^2)$. We know that the Gaussian is a positive definite function (its Fourier transform is a positive Gaussian).
Using the same logic as in (3): $\int \varphi(x) G(x) dx = \int \hat{G}(\xi) d\mu(\xi) \geq 0$ since $\hat{G}(\xi) = (\pi/\alpha)^{n/2} \exp(-|\xi|^2/4\alpha) > 0$ and $\mu$ is a positive measure.
\end{solution}

\begin{exercise}
For a $2\pi$-periodic function $x\in L^2([0,2\pi])$, let $x_n(t) = x(nt)$, $n=1,2,\ldots$.

Prove that $\{x_n\}$ converges weakly in $L^2([0,2\pi])$ and find the limit.
\end{exercise}

\begin{solution}
\textbf{Pedagogical Motivation:} This exercise illustrates the \textit{Riemann-Lebesgue Lemma} in the context of weak convergence. It shows that as a periodic function is compressed (increased frequency), its oscillation "averages out" when integrated against a fixed test function. This "averaging principle" is a fundamental concept in homogenization and multiscale analysis.

Let $x(t) \sim \sum_{k \in \mathbb{Z}} a_k e^{ikt}$ be the Fourier series of $x$, where $a_k = \frac{1}{2\pi} \int_0^{2\pi} x(t) e^{-ikt} dt$.
Then $x_n(t) = x(nt) \sim \sum_{k \in \mathbb{Z}} a_k e^{iknt}$.
To show weak convergence, we check the inner product $\langle x_n, y \rangle$ for any $y \in L^2([0, 2\pi])$.
Let $y(t) \sim \sum_{j \in \mathbb{Z}} b_j e^{ijt}$. By Parseval's identity:
\[ \langle x_n, y \rangle = \int_0^{2\pi} x_n(t) \overline{y(t)} dt = 2\pi \sum_{k \in \mathbb{Z}} a_k \overline{b_{kn}}. \]
As $n \to \infty$:
For $k = 0$, the term is $2\pi a_0 \overline{b_0}$, which is constant for all $n$.
For $k \neq 0$, the index $kn \to \pm \infty$. Since $y \in L^2$, we have $\sum |b_j|^2 < \infty$, which implies $\lim_{j \to \pm \infty} b_j = 0$ by the \textbf{Riemann-Lebesgue Lemma}.
Thus for each $k \neq 0$, $b_{kn} \to 0$ as $n \to \infty$.
Since $|a_k \overline{b_{kn}}| \leq \frac{1}{2}(|a_k|^2 + |b_{kn}|^2)$ and $\sum |a_k|^2 < \infty$, the sum $\sum_{k \neq 0} a_k \overline{b_{kn}}$ converges to 0 by the Dominated Convergence Theorem for sequences.
\[ \lim_{n \to \infty} \langle x_n, y \rangle = 2\pi a_0 \overline{b_0}. \]
Note that $2\pi a_0 = \int_0^{2\pi} x(t) dt$ and $2\pi \overline{b_0} = \int_0^{2\pi} \overline{y(t)} dt$.
The limit can be written as $\langle C, y \rangle$ where $C$ is the constant function:
\[ C = \frac{1}{2\pi} \int_0^{2\pi} x(t) dt. \]
Thus $x_n$ converges weakly to the constant function which is the average value of $x$.
\end{solution}

\begin{exercise}
Let $f$ be an entire function on $\mathbb{C}$ and there exists a constant $C>0$ such that $|f(z)|\leq C\sqrt{|z|}|\cos(z)|$ for all $z\in\mathbb{C}$. Prove that $f$ is identically zero.
\end{exercise}

\begin{solution}
\textbf{Pedagogical Motivation:} This problem combines the study of zeros of holomorphic functions with growth estimates (\textit{Liouville's Theorem}). By factoring out the zeros of $\cos(z)$, we transform a constrained growth condition into a standard polynomial growth condition, which forces the function to be a very simple polynomial (in this case, zero).

Consider the zeros of the function $\cos(z)$. The zeros are $z_k = (k + \frac{1}{2})\pi$ for $k \in \mathbb{Z}$.
At these points, the condition $|f(z_k)| \leq C \sqrt{|z_k|} |\cos(z_k)|$ implies $|f(z_k)| \leq 0$, so $f(z_k) = 0$.
Since $\cos(z)$ has only \textbf{simple zeros} (its derivative $-\sin(z)$ is non-zero at $z_k$), we can define a new function:
\[ g(z) = \frac{f(z)}{\cos(z)}. \]
Since every zero of $\cos(z)$ is also a zero of $f(z)$, $g(z)$ has only removable singularities at $z_k$. Thus $g(z)$ can be extended to an \textbf{entire function}.
The given condition implies:
\[ |g(z)| \leq C \sqrt{|z|} \text{ for all } z \in \mathbb{C}. \]
According to the \textbf{Generalized Liouville's Theorem}, if an entire function $g(z)$ satisfies $|g(z)| \leq M (1 + |z|^n)$, then $g(z)$ is a polynomial of degree at most $\lfloor n \rfloor$.
In our case, $n = 1/2 < 1$. Thus $g(z)$ must be a polynomial of degree at most 0, i.e., $g(z) = A$ for some constant $A$.
The condition $|A| \leq C \sqrt{|z|}$ must hold for all $z \in \mathbb{C}$.
Taking $z \to 0$ or considering the behavior as $z \to 0$, we get $A = 0$ (since $|A| \leq 0$ at $z=0$).
Thus $g(z) = 0$ for all $z$, which implies $f(z) = 0$ for all $z \in \mathbb{C}$.
\end{solution}

\begin{exercise}
Let $u(t,x)$ satisfy the following equation:
\[
u_t + \sum_{i=1}^n \psi_i(t,x,u)u_{x_i} = \mu\Delta u;\quad u(0,x) = u_0,
\]
where $u_0\in C^2(\mathbb{R}^n)$, $\psi_i$, $i=1,\ldots,n$ are of bounded $C^2$-norm, $\Delta = \sum_{i=1}^n \partial_{x_i}^2$ is the Laplacian in $\mathbb{R}^n$ and $\mu>0$. Assume $|u_0|\leq e^{\Phi/\mu}$ where $\Phi$ is bounded above and has Lipschitz constant $1$. Assume that $\left(\sum_{i=1}^n|\Psi_i(t,x,u)|^2\right)^{1/2}\leq A$, where $A$ is a positive constant.

Prove that there exists a constant $C$ depending only on $n$ such that
\[
|u(t,x)| \leq e^{Ct}e^{((A+1)t+\Phi(x)+2)/\mu}, \quad \forall t\geq 0.
\]
\end{exercise}

\begin{solution}
\textbf{Pedagogical Motivation:} This problem is a parabolic PDE estimate involving drift ($\psi$) and diffusion ($\mu \Delta$). It can be approached using the \textit{Maximum Principle} or the \textit{Feynman-Kac formula}. The goal is to show how the "amplitude" of the solution grows over time relative to the drift velocity $A$ and the diffusion coefficient $\mu$.

This problem can be solved using the comparison principle for parabolic equations or the Feynman-Kac representation. Let's use the comparison principle.
Let $v(t,x) = \mu \ln |u(t,x)|$ (formally). Then $u = \pm e^{v/\mu}$. Substituting into the equation $u_t + \psi \cdot \nabla u = \mu \Delta u$:
\[ v_t + \psi \cdot \nabla v = \mu \Delta v + \frac{1}{\mu} |\nabla v|^2. \]
At $t=0$, $v(0, x) = \mu \ln |u_0(x)| \leq \Phi(x)$.
Since $|\psi| \leq A$ and $|\nabla \Phi| \leq 1$, we expect the solution to shift by at most $At$ and spread by diffusion.
Alternatively, consider the heat kernel $K(t, x) = (4\pi \mu t)^{-n/2} \exp(-|x|^2 / 4\mu t)$.
The drift $\psi$ with $|\psi| \leq A$ implies that the mass can travel a distance at most $At$.
Since $\Phi$ is Lipschitz 1, $e^{\Phi(x-y)/\mu} \leq e^{(\Phi(x)+|y|)/\mu}$.
For the heat equation without drift ($w_t = \mu \Delta w$):
\[ w(t,x) = \int_{\mathbb{R}^n} K(t, y) e^{\Phi(x-y)/\mu} dy \leq e^{\Phi(x)/\mu} \int_{\mathbb{R}^n} K(t, y) e^{|y|/\mu} dy. \]
The integral $I = \int e^{-|y|^2/4\mu t + |y|/\mu} dy$ can be estimated. Completing the square in the exponent:
$- \frac{1}{4\mu t} (|y|^2 - 4t|y|) = - \frac{1}{4\mu t} (|y| - 2t)^2 + \frac{4t^2}{4\mu t} = - \frac{1}{4\mu t} (|y| - 2t)^2 + \frac{t}{\mu}$.
This gives a factor of $e^{t/\mu}$. The drift $A$ adds another factor $e^{At/\mu}$ because it effectively shifts the center of the kernel by $At$.
The $e^{Ct}$ term accounts for the integration in $n$ dimensions (the volume factor from the Jacobian and the polynomial terms in the Gaussian integral).
Specifically, $|u(t,x)| \leq (\text{Const}_n) e^{t/\mu} e^{At/\mu} e^{\Phi(x)/\mu}$.
For large $t$, any polynomial in $t$ can be bounded by $e^{Ct}$. The constant $+2$ in the provided exponent safely covers these normalization constants and polynomial factors for all $n$.
\end{solution}

